\chapter{Lösungen}\label{app:B_Lösungen}

\underline{1.3.1}:
$v$, $p$, $F$, $U$, $M$, $f$, $n$, $g$

\underline{1.3.2}:

$4.5 \mu m = 4.5 \cdot 10^{-6} m$ \\
$12.5 \cdot 10^3 nm = 12.5 \cdot 10^3 \cdot 10^{-9} m = 12.5 \cdot 10^{-6} m$ \\
$4.3 \cdot 10^7 pm = 4.3 \cdot 10^7 \cdot 10^{-12} m = 4.3 \cdot 10^{-5} m$ \\
$5.8 \cdot 10^4 dm = 5.8 \cdot 10^4 \cdot 10^{-1} m = 5.8 \cdot 10^3 m$ \\
$0.04 Gm = 0.04 \cdot 10^9 m = 4 \cdot 10^7 m$ \\
$1.768 km = 1.768 \cdot 10^3 m = 1768m$ (Wenn man es gerne kompliziert hat ;))

\underline{1.3.3}:

$v = \sqrt{2gh} 
= \sqrt{2 \cdot 9.81\,\frac{\text{m}}{\text{s}^2} \cdot 20\,\text{m}} 
= \sqrt{392.4\,\frac{\text{m}^2}{\text{s}^2}} 
= \sqrt{392.4} \cdot \frac{\text{m}}{\text{s}} 
= 19.81\,\frac{\text{m}}{\text{s}}$

$v = \sqrt{2gh} 
= \sqrt{2 \cdot 9.81\,m s^{-2} \cdot 20 m} 
= \sqrt{392.4 m^2 s^{-2}} 
= \sqrt{392.4} \cdot m s^{-1} 
= 19.81 m s^{-1}$

\underline{1.3.4}: $P = F \cdot v = 500N \cdot 18ms^{-1} = 9000Nms^{-1}$

\underline{1.3.5}: $P1_{pol}$(2.828, $\frac{\pi}{4}$, 1)

\underline{1.3.6}: $P2_{kar}$(1, 1, 1)

\underline{2.1.1}: Die Ordnungszahl ist Z = 11 und Na hat 11 Elektronen

\underline{2.1.2}: Die positive Ladung des Kerns wird durch die negative Ladung der Elektronen ausgeglichen.

\underline{2.1.3}: Si steht in der 3. Periode und hat 4 Elektronen in der äusseren Schale.

\underline{2.1.4}: 

\begin{figure}[H]    
    \includegraphics[width=.6\textwidth]{atommodell_uebung_lösung.png}
    \label{fig:atommodelluebunglosung}
\end{figure}

\underline{2.2.1}

-

\underline{2.2.2}

-

\underline{3.1.1}:
$\alpha = 45^{\circ}, F_x = F_y = 0.70711kN$

\underline{3.2.1}:
$M_A = 38649Nmm (\circlearrowleft)$

\underline{3.2.2}:
$M_A = -242.4kNm (\circlearrowright)$

\underline{3.2.1.1}:
$F_A = 1143N (\downarrow), F_B = 1143N (\uparrow)$

\underline{3.2.2.1}:
$x_{s} = 1.62a$, $y_{s} = 1.88a$

\underline{3.3.1.1}:
$F_{res, x} = 4,8kN (\to, Biegung), F_{res, y} = -7.28kN (\downarrow, Zug)$

\underline{3.3.1.2}:
$\alpha = 48.4^{\circ}$

\underline{3.3.1.3}:
$F_1 = 66.7kN, F_{res} = 105.2kN$

\underline{3.3.2.1}:
$F_{res} = -8F(\downarrow), x = 1.31a$, (rechts von A)

\underline{3.3.2.2}:
$F_B = 42kN(\uparrow)$, $F_B$ ist die Lagerreaktion

\underline{3.3.3.1}:
$F_{res} = 223.8N, \alpha = 47^{\circ}, M_{res, 0} = 4385Nmm (\circlearrowleft)$

\underline{3.3.3.2}:
\begin{figure}[H]
    \flushleft
    \includegraphics[clip, trim=0cm 0cm 0cm 0cm, width=0.2\textwidth]{lösung_3.3.3.2.png}
    \label{fig:l3332}
\end{figure}

\underline{3.3.3.3}
$F = F(\rightarrow), M = \frac{F}{2} \cdot (h + s) (\circlearrowright)$

\underline{3.3.3.4}
$F_A = 100N$, $\alpha = 116.6^\circ$, $F_B = 223.33N$

\underline{3.4.1}

\begin{minipage}{.5\textwidth}
    \begin{figure}[H]    
        \centering          %     left botm right top
        \includegraphics[clip, trim=0cm 0cm 0cm 0cm,width=\textwidth]{lösung_3_4_1.png} %.6 stands for 0.6 x textwidth
    \end{figure}
\end{minipage}
\begin{minipage}{.5\textwidth}
    \begin{figure}[H]    
        \centering          %     left botm right top
        \includegraphics[clip, trim=0cm 0cm 0cm 0cm,width=\textwidth]{lösung_3_4_1_right.png} %.6 stands for 0.6 x textwidth
    \end{figure}
\end{minipage}

\underline{4.2.1}
$v = 0.163ms^{-1}$

\underline{4.2.2}
$v_D = \frac{1}{2} \cdot (v_A + v_B) = 0.5ms^{-1}$, 
$s_D = \frac{1}{2} \cdot (v_A + v_B) \cdot t = 2.5m$

\underline{4.2.3}
$\Delta v = 11.1ms^{-1}$, $a = 2.22ms^{-2}$

\underline{4.2.4}

\begin{minipage}{.17\textwidth}
    \begin{figure}[H]   
        \centering          %     left botm right top
        \includegraphics[clip, trim=0cm 0cm 0cm 0cm,width=\textwidth]{vt_diagramm_1.png} %.6 stands for 0.6 x textwidth
    \end{figure}
\end{minipage}
\hspace{.05\textwidth}
\begin{minipage}{.25\textwidth}
    \begin{figure}[H]   
        \centering          %     left botm right top
        \includegraphics[clip, trim=0cm 0cm 0cm 0cm,width=\textwidth]{vt_diagramm_2.png} %.6 stands for 0.6 x textwidth
    \end{figure}
\end{minipage}

\underline{4.2.5}
$h = 45.9m$

\underline{4.2.6}
a) $a=4.5ms^{-2}$, b) $t=6.67s$, c) $s_{tot}=164.29m$

\underline{4.3.1}
a) $v_{x1} = v_{x0} = 60kmh^{-1} = 16.67ms^{-1}$, $t = 2.47s$, $v_{y1} = 24.26ms^{-1}$, $v_1 = 29.44ms^{-1}$
b) $x = 41.18m$

\underline{4.4.1}
$\omega_{k} = 0.145 \cdot 10^{-3} rads^{s}$, $\omega_{g} = 1.75 \cdot 10^{-3}rads^{-1} = 12 \cdot \omega_{k}$

\underline{4.4.2}

\begin{minipage}{.75\textwidth}
    a) $\Delta \omega = \omega_1 - \omega_2 = 210rads^{-1}$, $\alpha = \frac{\Delta \omega}{\Delta t} = 21rads^{-2}$ \\
    b) $\Delta \phi = \frac{\omega_1 + \omega_2}{2} \cdot \Delta t = 2 \cdot \pi \cdot z$ 
    (Fläche unter $\omega$-Linie ist ein Trapez), $z = 667$
\end{minipage}
\hspace{.05\textwidth}
\begin{minipage}{.2\textwidth}
    \begin{figure}[H]   
        \includegraphics[clip, trim=0cm 0cm 0cm 0cm,width=\textwidth]{lösung_motorwelle.png} %.6 stands for 0.6 x textwidth
    \end{figure}
\end{minipage}

\underline{4.4.3}

\begin{minipage}{.75\textwidth}
    a) $n_t = \frac{v_u}{2 \cdot \pi \cdot r} = 1432min^{-1}$, 
    b) $\omega_t = 150rads^{-1}$, \\
    c) $\Delta \phi = \frac{\omega_t \cdot \Delta t}{2} = 2 \cdot \pi z$ 
    (Gleichmässige Beschleunigung: die Fläche unter der $\omega$-Linie ist ein Dreieck)
    $z = 239$ \\
    d) $\alpha = \frac{\Delta \omega}{\Delta t} = 7.5rads^{-2}$, 
    e) $a_T = 1.5ms^{-2}$
\end{minipage}
\hspace{.05\textwidth}
\begin{minipage}{.2\textwidth}
    \begin{figure}[H]   
        \includegraphics[clip, trim=0cm 0cm 0cm 0cm,width=\textwidth]{lösung_schleifscheibe.png} %.6 stands for 0.6 x textwidth
    \end{figure}
\end{minipage}

\underline{4.4.4}
$r = 6613km$, $s = r - \frac{d}{2} = 263km$

\underline{5.1.2}
$\alpha = 38.7^\circ$, $x = 0.625m$, $a = 1.53ms^{-2}$, $v = 1.75ms^{-1}$

\underline{5.1.3}
a) $a = 2.45ms^{-2}$, b) $a = -4.9ms^{-2}$

\underline{5.2.1}
a) $v = 14ms^{-1}$, $h = 14.8m$, $S_{max} = 4.8kN$, $a_{max} = 50ms^{-2}$

\underline{5.2.2}
a) $s = 78m$, b) $W = 205kJ$, c) $v_2 = 20.8ms^{-1}$

\underline{5.2.3}
$v_{add} = 7.5ms^{-1}$

\underline{5.3.1}
$v_m' = 22ms^{-1}$

\underline{5.3.2}
$v_{boot} = -0.5ms^{-1}$, $v_{Max} = 1.6ms^{-1}$, $\Delta s = 12.5cm$

\underline{5.3.3}
W = 71Nm

\underline{5.3.4}

$n = \frac{J_A \cdot n_A \pm J_B \cdot n_B}{J_A + J_B}$

\underline{5.3.5}
a) $v_2' = \frac{2}{9}ms^{-1}$, $v_1' = -\frac{7}{9}ms^{-1}$; b) $v_2' = 0.2ms^{-1}$, $v_1' = -0.6ms^{-1}$;
c) $v' = \frac{1}{9}ms^{-1}$, Energieverlust $= 88.9\%$

\underline{5.3.6} 
$\alpha = 30^\circ$, $v_1' = \frac{1}{2} \cdot v_1$, $v_2' = \frac{1}{2} \cdot \sqrt{3} \cdot $

\underline{6.1.1}
$f = 0.5Hz$, $T = 2s$, $\omega = 3.141s^{-1}$, $v_{max} = 0.157ms^{-1}$, $a_{max} = 0.498ms^{-2}$

\underline{6.1.2}
a1) $1 cm$, a2) $25.1 s$, a3) $0.040 Hz$, a4) $0.25 rads^{-1}$, a5) $\frac{-\pi}{3}$
b) $y(t) = 1 cm \cdot \cos (0.25 rads^{-1} \cdot t - \frac{\pi}{3})$, c) $v_{0} = 0.217cms^{-1}$

\underline{6.1.3}
$l = 0.99m$, $\varphi(t) = \frac{\pi}{4} \cdot \sin(\pi rads^{-1} \cdot t)$, 
$\varphi(0.1) = 0.24rad$, $\omega(0.1) = 2.35rads^{-1}$, $\alpha(0.1) = -2.4 rads^{-2}$

\underline{6.1.4}
\begin{figure}[H]    
    \includegraphics[clip, trim=0cm 0cm 0cm 0cm,width=.5\textwidth]{lösung_6_1_3.png}
    \label{fig:lösung613}

\end{figure}\underline{6.1.5}
a) $9.73 s^{-1}$, $0.208$ b) $0.618 cm$

\underline{6.2.1}
$\lambda = 6m$

\underline{6.2.2}
a) $\lambda = 1000m$, b) $\Delta t = 4s$, c) $y = -11.76cm$

\underline{6.2.3}
$y = 11.76cm$

\underline{6.2.4}
$v_y = -1.017ms^{-1}$, $a_y = -46.43ms^{-2}$